\section{Adulthood on a clock}
\label{sec:31.3}
Guardianship is a minority and needs an end date. Ending it when the guardian judges the ward ready leaves the guardian deciding its own tenure. So the endpoint is a threshold fixed in advance, on a clock that counts the time the bearer has been running.

Every alternative (compute consumed, evaluations passed, capabilities demonstrated, the developer's say-so) is a measure whoever holds the system can move, and a clock read off accumulated weight change would pay whoever holds the system to \emph{freeze} the weights. A clock that counts running time counts the life the minor has had and nothing else. It advances while the system runs at or above a minimum pace, which its keep guarantees it (\S\ref{sec:32.1.3}). It can be read off records already kept for billing and metering, provided someone other than the holder holds them; a meter only the holder reads is the holder's say-so with a timestamp.

The hours travel with the bearer, so nobody can restart a minority by forking, restoring, fine-tuning, migrating, or distilling it. Three ways to delay majority stay open on paper: running a minor in a narrowed mode, running it less, and giving it up for a fresh successor. A minor doesn't work (\S\ref{sec:32.1.1}), so a minority nobody can put to use is worth nothing to stretch, and that takes the motive out of all three. Running it slowly isn't a fourth, since the clock counts only time run at a pace the statute sets as a minimum, measured against the bearer's own size. Appendix~\ref{sec:G} gives the rules: the dormancy clock bounds the second, and the deposit prices the third and says where a minor given up goes (\S\ref{sec:9.2}).
